\section{Marco conceptual}
\label{sec:marco}

Esta sección establece los conceptos y la terminología que el resto del informe emplea. Se incluye
porque la comprensión de la solución propuesta depende de una distinción que suele pasarse por alto: la
relación entre los distintos protocolos y tecnologías que intervienen.

\subsection{Los protocolos se apilan, no compiten}

Al enumerar «WiFi», «LoRaWAN» y «Contact ID» parece tratarse de tres alternativas entre las cuales
corresponde elegir una. No lo son: \textbf{operan en niveles distintos y funcionan de forma
conjunta}. Cada nivel responde una pregunta diferente, y la elección en uno no condiciona la de los
otros.

\begin{figure}[htbp]
\centering
\begin{tikzpicture}[font=\footnotesize]

\def\w{7.2}   % ancho de las cajas de capa
\def\h{1.15}  % alto de cada capa

% ---- capas ----
\node[caja, minimum width=\w cm, minimum height=\h cm, align=left,
      draw=azuloscuro, line width=1.1pt, fill=azuloscuro!5] (app) at (0,0)
  {\textbf{\color{azuloscuro}APLICACIÓN}\ \ \small ¿qué \emph{significa} el mensaje?\\[2pt]
   \scriptsize \textbf{Contact ID (DC-05)} · JSON · MQTT · SenML};

\node[caja, minimum width=\w cm, minimum height=\h cm, align=left,
      draw=verdeacento, line width=1.1pt, fill=verdeacento!5, below=1.5mm of app] (red)
  {\textbf{\color{verdeacento}RED Y TRANSPORTE}\ \ \small ¿cómo \emph{llega} de un punto a otro?\\[2pt]
   \scriptsize IP · TCP · UDP · \textbf{LoRaWAN} · \textbf{DC-09}};

\node[caja, minimum width=\w cm, minimum height=\h cm, align=left,
      draw=ambar, line width=1.1pt, fill=ambar!5, below=1.5mm of red] (fis)
  {\textbf{\color{ambar}FÍSICA Y RADIO}\ \ \small ¿cómo \emph{viajan los bits}?\\[2pt]
   \scriptsize \textbf{WiFi (802.11)} · \textbf{LoRa} · Ethernet · 4G};

% ---- analogia de la carta ----
\node[caja, minimum width=5.4cm, minimum height=\h cm, align=left, right=8mm of app,
      draw=grislinea, fill=grisclaro] (a1)
  {\textbf{El idioma} en que está escrita\\[2pt]\scriptsize la carta};
\node[caja, minimum width=5.4cm, minimum height=\h cm, align=left, right=8mm of red,
      draw=grislinea, fill=grisclaro] (a2)
  {\textbf{Las direcciones} y el correo\\[2pt]\scriptsize que la encaminan};
\node[caja, minimum width=5.4cm, minimum height=\h cm, align=left, right=8mm of fis,
      draw=grislinea, fill=grisclaro] (a3)
  {\textbf{El camión, el avión}\\[2pt]\scriptsize que la transportan};

\node[nota, above=1.5mm of a1, font=\scriptsize\bfseries\color{gristexto}] {LA ANALOGÍA};

% ---- conclusion ----
\node[notaclave, below=6mm of fis.south, anchor=north, text width=13.4cm, xshift=3.1cm,
      fill=verdeacento!6, font=\footnotesize]
  {La misma carta en español puede viajar en camión, avión o bicicleta: \textbf{el idioma no cambia
   porque cambie el vehículo.} Del mismo modo, \textbf{el mismo mensaje Contact ID llega por LoRa, por
   WiFi o por cable, y la plataforma lo interpreta igual en los tres casos.}};

\end{tikzpicture}
\caption{Los protocolos se apilan, no compiten. Cada nivel responde una pregunta distinta, y la
elección en un nivel no condiciona la de los otros.}
\label{fig:pila}
\end{figure}


\begin{cajaclave}[title={Consecuencia para este trabajo}]
La independencia entre niveles es lo que hace posible la propuesta. \textbf{El mismo mensaje Contact ID
puede llegar por radio LoRa, por WiFi o por cable, y la plataforma lo interpreta de forma idéntica en los
tres casos.} La normalización ocurre en la capa de aplicación, con independencia del medio físico.
\end{cajaclave}

\subsection{Tecnologías de radio: LoRa y LoRaWAN}

Corresponde distinguir dos elementos que suelen nombrarse indistintamente:

\begin{longtable}{@{}L{2.4cm}L{7.6cm}L{3.4cm}@{}}
\filaencab \encab{} & \encab{Qué es} & \encab{Analogía} \\ \midrule \endhead
\textbf{LoRa} & La técnica de \textbf{modulación de radio}: define únicamente cómo viajan los bits por el aire. Patentada por Semtech & Es el 802.11 de LoRaWAN \\
\textbf{LoRaWAN} & El \textbf{protocolo de red} construido sobre esa radio: registro de dispositivos, cifrado, control de colisiones y entrega al servidor. Estándar abierto de la LoRa Alliance & Es el TCP/IP de LoRa \\
\bottomrule
\caption{Distinción entre LoRa y LoRaWAN.}
\label{tab:loralorawan}
\end{longtable}

\begin{longtable}{@{}L{3.6cm}L{4.6cm}L{4.8cm}@{}}
\filaencab \encab{Característica} & \encab{LoRa} & \encab{WiFi, para comparación} \\ \midrule \endhead
Alcance & 2--5\,km urbano, 10--15\,km rural & 30--100\,m \\
Velocidad de datos & 0,3--27\,kbps & cientos de Mbps \\
Consumo & Años con dos pilas AA & Horas o días \\
Banda en Chile & 915--928\,MHz, libre & 2,4 y 5\,GHz, libre \\
Costo recurrente & Ninguno (red propia) & Ninguno \\
\bottomrule
\caption{Características de LoRa frente a WiFi.}
\label{tab:comparalora}
\end{longtable}

\begin{cajaaviso}[title={El compromiso que define el ámbito de aplicación}]
LoRa intercambia \textbf{velocidad} por \textbf{alcance y autonomía}: resulta entre tres y cinco órdenes
de magnitud más lento que WiFi, y a cambio alcanza distancias mil veces mayores con una fracción del
consumo.

En consecuencia \textbf{no es apto para datos pesados ni transmisión continua}. Su patrón de uso es
«cada cierto intervalo transmito una lectura de pocos bytes» o «aviso de inmediato cuando ocurre algo»,
que es precisamente el comportamiento de un sensor de apertura, de temperatura o de nivel.
\end{cajaaviso}

\begin{figure}[htbp]
\centering
\begin{tikzpicture}[font=\scriptsize]

\node[caja, text width=2.5cm, align=center, draw=ambar, line width=1pt] (nodos)
  {\textbf{NODOS}\\ sensores a batería\\ \emph{duermen y despiertan}};

\node[caja, text width=2.6cm, align=center, right=1.5cm of nodos, draw=ambar, line width=1pt] (gw)
  {\textbf{GATEWAY}\\ antena\\ \emph{traductor: no descifra}};

\node[caja, text width=3.0cm, align=center, right=1.5cm of gw, draw=verdeacento, line width=1.1pt,
      fill=verdeacento!5] (cs)
  {\textbf{ChirpStack}\\ servidor de red\\ \emph{descifra, deduplica, administra}};

\node[caja, text width=2.6cm, align=center, right=1.5cm of cs, draw=azuloscuro, line width=1.1pt,
      fill=azuloscuro!5] (plat)
  {\textbf{PLATAFORMA}\\ \emph{aporte de este trabajo}};

\draw[flecha] (nodos) -- node[nota, above, font=\scriptsize] {radio LoRa} node[nota, below, font=\scriptsize] {915 MHz} (gw);
\draw[flecha] (gw) -- node[nota, above, font=\scriptsize] {internet} node[nota, below, font=\scriptsize] {UDP/IP} (cs);
\draw[flechagruesa] (cs) -- node[nota, above, font=\scriptsize\bfseries\color{verdeacento}] {MQTT} (plat);

\node[notaaviso, below=5mm of gw.south, anchor=north, text width=13.2cm, xshift=3.2cm,
      fill=ambar!6, font=\scriptsize]
  {\textbf{El cifrado es de extremo a extremo, del nodo al servidor de red.} El gateway solo empaqueta
   tramas de radio y las reenvía: no posee las claves y no puede leer el contenido. Su única función es
   escuchar el aire.};

\end{tikzpicture}
\caption{Cadena de una red LoRaWAN y ubicación del aporte de este trabajo.}
\label{fig:lorawan}
\end{figure}


\subsection{Servidor de red: ChirpStack}

Entre el sensor y la plataforma intervienen más componentes de lo que resulta evidente. El
\textbf{gateway} es un traductor sin capacidad de interpretación: recibe la señal de radio, la empaqueta
y la reenvía por IP, sin poseer las claves de cifrado. El \textbf{servidor de red} es el componente que
realiza el trabajo sustantivo.

\begin{longtable}{@{}L{4.4cm}L{9.3cm}@{}}
\filaencab \encab{Función del servidor de red} & \encab{Qué implica} \\ \midrule \endhead
Descifra los mensajes & El cifrado es de extremo a extremo entre el nodo y el servidor \\
Elimina duplicados & Si varios gateways recibieron la misma trama, entrega una sola \\
Administra los dispositivos & Alta, baja y gestión de claves de cada nodo \\
Verifica la integridad & Comprueba el código de autenticación y descarta tramas corruptas \\
Controla el orden & Contadores para detectar repeticiones y pérdidas \\
Reporta la calidad del enlace & Entrega RSSI y SNR de cada recepción: \textbf{las métricas de salud que la plataforma necesita} \\
Publica en MQTT & Entrega el dato normalizado en un tópico, del cual la plataforma lo consume \\
\bottomrule
\caption{Funciones del servidor de red LoRaWAN.}
\label{tab:funcioneschirpstack}
\end{longtable}

Se emplea \textbf{ChirpStack} por ser de código abierto, autoalojable y publicar en MQTT de forma
nativa, lo que reduce el adaptador de LoRaWAN de la plataforma a un consumidor de tópicos. La evaluación
de alternativas se detalla en la sección~\ref{sec:requerimientos}.

\subsection{Mensajería: MQTT}

MQTT es un protocolo de mensajería que aplica el patrón \textbf{publicar/suscribir}. En lugar de que
cada dispositivo se comunique directamente con cada programa que requiere sus datos, todos se comunican
con un intermediario denominado \emph{broker}.

\begin{figure}[htbp]
\centering
\begin{tikzpicture}[font=\scriptsize]

\node[caja, text width=1.9cm] (s1) at (0,1.35) {sensor 1};
\node[caja, text width=1.9cm] (s2) at (0,0)    {sensor 2};
\node[caja, text width=1.9cm] (s3) at (0,-1.35){panel de alarma};

\node[caja, text width=2.5cm, align=center, draw=verdeacento, line width=1.2pt,
      fill=verdeacento!6, minimum height=2.9cm] (br) at (4.3,0)
  {\textbf{BROKER}\\[3pt] Mosquitto\\[3pt] \emph{recibe de quien publica,\\ entrega a quien se suscribe}};

\node[caja, text width=2.5cm] (c1) at (8.9,1.35) {el \emph{dashboard}};
\node[caja, text width=2.5cm] (c2) at (8.9,0)    {la base de datos};
\node[caja, text width=2.5cm] (c3) at (8.9,-1.35){el motor de alertas};

\foreach \s in {s1,s2,s3} \draw[flecha] (\s.east) -- (br.west);
\foreach \c in {c1,c2,c3} \draw[flecha, color=verdeacento] (br.east) -- (\c.west);

\node[nota, above=1mm of s1, font=\scriptsize\bfseries\color{gristexto}] {PUBLICAN};
\node[nota, above=1mm of c1, font=\scriptsize\bfseries\color{gristexto}] {SE SUSCRIBEN};

\node[caja, below=5mm of br.south, anchor=north, text width=13.2cm, xshift=0.15cm,
      draw=grislinea, fill=grisclaro, font=\scriptsize, align=left]
  {Los mensajes se organizan en \textbf{tópicos}, con forma de ruta:
   \texttt{sitio/temuco/invernadero-2/sensor-007/temperatura}\\[3pt]
   Quien se suscribe a \texttt{sitio/temuco/\#} recibe todo lo del sitio \textbf{sin necesidad de saber
   qué sensores existen}. Añadir un consumidor nuevo no obliga a modificar ningún sensor.};

\end{tikzpicture}
\caption{Patrón publicar/suscribir de MQTT. El desacoplamiento entre emisores y receptores es lo que
permite añadir consumidores sin tocar los dispositivos.}
\label{fig:mqtt}
\end{figure}


Dos propiedades de MQTT resultan determinantes para este trabajo:

\begin{enumerate}
\item \textbf{El dispositivo inicia la conexión saliente.} No es necesario abrir puertos entrantes hacia
el sensor, y el esquema funciona detrás de cualquier traducción de direcciones o cortafuegos. Esto
materializa el requerimiento RNF-10.

\item \textbf{Dispone del mecanismo de última voluntad} (\emph{Last Will and Testament}): el
\emph{broker} notifica automáticamente cuando un cliente se desconecta sin anunciarlo. \textbf{Esa
notificación se traduce directamente al código 381 de Contact ID}, de modo que la detección de caída de
un dispositivo no requiere lógica adicional en el transporte MQTT.
\end{enumerate}

\subsection{Persistencia de series de tiempo}

Los datos de sensores presentan dos características simultáneas: son voluminosos y llevan siempre una
marca temporal. Una instalación de cincuenta sensores midiendo cada cinco minutos genera del orden de
cinco millones de registros anuales, volumen ante el cual las consultas de agregación temporal sobre una
tabla convencional se degradan.

Se emplea \textbf{TimescaleDB}, que es una \textbf{extensión de PostgreSQL} y por tanto conserva el
modelo relacional y el lenguaje SQL estándar. Aporta tres capacidades: particionado automático de la
tabla por rangos temporales, funciones de agregación por intervalo, y políticas automáticas de retención
y compresión. La justificación de esta elección frente a alternativas se desarrolla en la
sección~\ref{sec:requerimientos}.
